\subsubsection{Registro de Análisis Pureza}

\noindent\textbf{Caso de uso:} Cargar Resultados de Pureza

\vspace{0.5em}
\noindent\textbf{Descripción:} El analista registra pesos y componentes de una muestra de semillas para calcular el porcentaje de pureza física, identificando semilla pura, materia inerte, otros cultivos y malezas.

\vspace{0.5em}
\noindent\textbf{PreCondición:}
\begin{itemize}[nosep]
    \item Debe existir al menos un lote activo con tipo de análisis ``PUREZA'' asignado
    \item El lote no debe tener un análisis de Pureza en estado APROBADO
    \item El usuario debe estar autenticado como Analista o Administrador
\end{itemize}

\vspace{0.5em}
\noindent\textbf{PostCondición:}
\begin{itemize}[nosep]
    \item Análisis creado con estado \texttt{EN\_PROCESO}
    \item Fecha de inicio automática asignada
    \item Resultados guardados y vinculados al lote
    \item Registro automático en historial de análisis
\end{itemize}

\vspace{0.5em}
\noindent\textbf{Actores:} Analista, Administrador.

\vspace{0.5em}
\noindent\textbf{Flujo de evento principal:}
\begin{enumerate}[nosep]
    \item \textbf{Fase 1: Creación del análisis}
    \begin{itemize}[nosep]
        \item El usuario selecciona un lote elegible para Pureza
        \item Accede a ``Crear análisis de Pureza''
        \item El sistema crea el análisis con estado \texttt{EN\_PROCESO} y fecha de inicio automática
    \end{itemize}
    \item \textbf{Fase 2: Registro de datos INIA}
    \begin{itemize}[nosep]
        \item Ingresa fecha del análisis (obligatoria)
        \item Ingresa peso inicial en gramos (obligatorio, debe ser > 0)
        \item Ingresa componentes en gramos: Semilla pura, Materia inerte, Otros cultivos, Malezas, Malezas toleradas, Malezas tolerancia cero
        \item El sistema calcula automáticamente: peso total y porcentajes sin redondeo (4 decimales)
        \item Validaciones en tiempo real: peso total $\leq$ peso inicial; alerta si pérdida de masa > 5\% (informativa)
        \item El usuario ingresa porcentajes redondeados manualmente; el sistema valida suma = 100\%
        \item Valores entre 0 y 0.05\% muestran ``tr'' (trazas) y son de solo lectura
    \end{itemize}
    \item \textbf{Fase 3: Datos INASE (opcionales)}
    \begin{itemize}[nosep]
        \item El usuario puede ingresar fecha INASE y porcentajes oficiales INASE
    \end{itemize}
    \item \textbf{Fase 4: Registro de listados (opcional)}
    \begin{itemize}[nosep]
        \item Registro de malezas identificadas y otros cultivos, indicando catálogo/especie, cantidad y categorías
    \end{itemize}
    \item \textbf{Fase 5: Cumplimiento del estándar}
    \begin{itemize}[nosep]
        \item El usuario selecciona si cumple o no con el estándar y el sistema guarda los cambios
    \end{itemize}
\end{enumerate}

\vspace{0.5em}
\noindent\textbf{Flujos alternativos:}
\begin{itemize}[nosep]
    \item \textit{Peso total mayor que peso inicial}: bloqueo de guardado y mensaje de error
    \item \textit{Pérdida de masa excesiva}: muestra alerta informativa (no bloquea)
    \item \textit{Suma de porcentajes redondeados $\neq$ 100\%}: muestra alerta en rojo (no bloquea)
    \item \textit{Finalizar sin evidencia}: sistema rechaza la finalización si faltan datos
\end{itemize}

\vspace{0.5em}
\noindent\textbf{Requerimientos especiales:}
\begin{itemize}[nosep]
    \item Cálculos automáticos en tiempo real con 4 decimales
    \item Validación de coherencia en pesos (total $\leq$ inicial)
    \item Alerta de pérdida de masa > 5\% y manejo de ``tr'' para trazas
    \item Tabla de tolerancias disponible para consulta
\end{itemize}
